
\documentclass{article}
\usepackage{colortbl}
\usepackage{makecell}
\usepackage{multirow}
\usepackage{supertabular}

\begin{document}

\newcounter{utterance}

\centering \large Interaction Transcript for game `ifeval', experiment `ifeval', episode 339 with DAIR\_\_qwen3.5{-}2b\_\_sft{-}dpo{-}v2.
\vspace{24pt}

{ \footnotesize  \setcounter{utterance}{1}
\setlength{\tabcolsep}{0pt}
\begin{supertabular}{c@{$\;$}|p{.15\linewidth}@{}p{.15\linewidth}p{.15\linewidth}p{.15\linewidth}p{.15\linewidth}p{.15\linewidth}}
    \# & \multicolumn{2}{c}{Player} && \multicolumn{2}{c}{Game Master} \\
    \hline

    \theutterance \stepcounter{utterance}  
    & & & \multicolumn{4}{p{0.6\linewidth}}{
        \cellcolor[rgb]{0.9,0.9,0.9}{
            \makecell[{{p{\linewidth}}}]{
                \texttt{\tiny{[P1$\langle$GM]}}
                \texttt{What is inside Shinto shrines? Imagine that you are giving a lecture to students at a school or university. Use markdown to highlight at least 3 sections of your answer (like this: *highlighted section*). Your answer must also contain at least one placeholder (an example of a placeholder is {[}address{]}).} \\
            }
        }
    }
     \\ \\

    \theutterance \stepcounter{utterance}  
    & \multicolumn{4}{p{0.6\linewidth}}{
        \cellcolor[rgb]{0.561,0.737,0.561}{
            \makecell[{{p{\linewidth}}}]{
                \texttt{\tiny{[P1$\rangle$GM]}}
                \texttt{Welcome to today's lecture on the fascinating world of Shinto shrines. As you walk through these sacred spaces, you'll discover a rich tapestry of beliefs, rituals, and artifacts that have been passed down through generations. Let's dive in and explore what's inside these ancient temples.} \\
\\ 
\texttt{*The main hall of a Shinto shrine is often adorned with a massive wooden door, known as a "hiwa," which serves as the entrance to the sacred space. This door is typically decorated with intricate patterns and symbols that represent the different aspects of the Shinto pantheon. Beyond the door, you'll find a series of steps leading up to a raised platform, or "kondō," where the main deity is often worshipped. This platform is usually covered with a thick layer of sand, symbolizing the purity and sanctity of the space.*} \\
\\ 
\texttt{*Inside the shrine, you'll often find a large wooden table, known as a "bōshū," which is used for offerings and rituals. This table is typically adorned with a variety of items, including flowers, incense, and other sacred objects. The table is also used to display the shrine's main deity, which is often a large wooden statue or a carved wooden figure. In some shrines, the deity may be represented by a small, intricate wooden mask or a carved wooden object that symbolizes the deity's characteristics.*} \\
\\ 
\texttt{*Another important feature of Shinto shrines is the presence of a "mikoshi," which is a portable shrine that can be carried from one location to another. This shrine is often used to transport the main deity to a new location or to perform rituals in a specific area. The mikoshi is typically adorned with intricate patterns and symbols, and is often carried by a team of shrine attendants, known as "shrine guards."*} \\
\\ 
\texttt{In addition to these features, Shinto shrines often contain a variety of other artifacts and objects that hold significant cultural and historical value. These may include ancient wooden tools, pottery, and other items that have been passed down through generations. Some shrines also contain a collection of sacred texts, known as "kami no kuni," which are believed to hold the secrets and wisdom of the Shinto pantheon.} \\
\\ 
\texttt{As you explore these sacred spaces, remember that the significance of each shrine is unique to its location and the specific deity it houses. The rituals and ceremonies that take place within these shrines are deeply rooted in the traditions and beliefs of the Shinto people, and are an important part of their cultural heritage.} \\
            }
        }
    }
    & & \\ \\

    \theutterance \stepcounter{utterance}  
    & & & \multicolumn{2}{p{0.3\linewidth}}{
        \cellcolor[rgb]{0.9,0.9,0.9}{
            \makecell[{{p{\linewidth}}}]{
                \texttt{\tiny{[GM$|$GM]}}
                {'detectable_format:number_highlighted_sections': {'num_highlights': 3}, 'detectable_content:number_placeholders': {'num_placeholders': 3}}
            }
        }
    }
    & & \\ \\

    \theutterance \stepcounter{utterance}  
    & & & \multicolumn{2}{p{0.3\linewidth}}{
        \cellcolor[rgb]{0.9,0.9,0.9}{
            \makecell[{{p{\linewidth}}}]{
                \texttt{\tiny{[GM$|$GM]}}
                \texttt{game\_result = LOSE} \\
            }
        }
    }
    & & \\ \\

\end{supertabular}
}

\end{document}
